% Source-aligned English for printed p.199.
\clearpage
\SourcePage{199}

\LayerTarget{tohoku:construction:canonical-map-inverse-direct:p1}
$g(x)$ on $V$, whose value at each $x\in V$ is an element of $f_*^G(A)(f(x))=\mathop{\underset{\to}{\mathrm{lim.}}}\limits_{U\ni f(x)}\Gamma(f^{-1}(U),A)^G$,
and such that for every $x\in V$ there exist an open neighborhood
$V'\subset V$ of $x$ and a section $h$ of $f_*^G(A)$ over a neighborhood
$U'$ of $f(x)$ (i.e. an invariant section $h$ of $A$ over $f^{-1}(U')$) such
that $g(x)=h(x)$ for $x\in V'\cap f^{-1}(U')$. We obtain a canonical
monomorphism\textsuperscript{10)}:

\LayerTarget{tohoku:equation:5.1.2}
\[
f^{-1}(f_*^G(A))\to A
\tag{5.1.2}
\]
\LayerTarget{tohoku:criterion:inverse-direct-essential-image}
This identifies $f^{-1}(f_*^G(A))$ with a subsheaf $A'$ of $A$. Formula
(5.1.1) above implies that $A'=A$ if and only if $A$ is isomorphic to a
$G$-sheaf of the form $f^{-1}(B)$. If this is so, then for every $x\in X$ the
operations of $G_x$ on $A(x)$ are trivial. The converse is also true if $G$
satisfies condition (D) below (cf. 5.3), as is readily checked;
\LayerTarget{tohoku:result:fixed-point-free-equivalence}
\textit{thus, if in addition every $G_x$ is reduced to the identity} (i.e. if
no $g\ne e$ in $G$ has a fixed point), \textit{the functors $f_*^G$ and
$f^{-1}$ establish mutually inverse isomorphisms between the category of
$G$-sheaves on $X$ and the category of sheaves on $Y$.}

\LayerTarget{tohoku:example:constant-direct-image}
\textsc{Examples of direct images of G-sheaves.} If $A$ is the
constant sheaf on $X$ defined by a set $M$ on which $G$ acts trivially, then
$A^G$ is the constant sheaf on $Y$ defined by the same set $M$. But if $G$
does not act trivially on $M$, and if the action on $X$ is not fixed-point
free, then $f_*^G(M)$ is in general no longer locally constant.
\LayerTarget{tohoku:example:germ-function-direct-images}
Suppose condition (D) of 5.3 holds, and let $A$ be the sheaf of germs of maps
from $X$ to a set $M$; then $A^G$ is the sheaf of germs of maps from $Y$ to
$M$. Under this correspondence, if $X$ is a differentiable manifold (resp. a
real analytic space, resp. a complex analytic space), and if $A$ is the sheaf
of germs of differentiable functions (resp. real analytic functions, resp.
complex analytic functions) with values in a space of the same kind, then
$A^G$ is the corresponding sheaf on $Y$. This last example is particularly
important, and the analogous statement is true (by definition, as for the
preceding examples) in abstract Algebraic Geometry, and even for arithmetic
varieties.

\LayerTarget{tohoku:section:5.2}
\bigskip
\textbf{5.2. The functors $H^n(X;G,A)$ and $\mathbf H^n(G,A)$ and the
fundamental spectral sequences.}
\LayerTarget{tohoku:definition:section-invariant-functors}
To avoid confusion, we write $\Gamma_X$ and $\Gamma_Y$ for the ``section''
functors defined for sheaves on $X$ and on $Y$, respectively, and $\Gamma^G$
for the functor $M\to M^G$ that associates to a set $M$ on which $G$
acts its subset of invariants. Finally, if $A$ is a $G$-sheaf on $X$, set

\LayerTarget{tohoku:equation:5.2.1}
\[
\Gamma_X^G(A)=\Gamma_X(A)^G
\tag{5.2.1}
\]
so that, by definition, we have the formulas

\LayerTarget{tohoku:equation:5.2.2}
\[
\Gamma_X^G=\Gamma^G\Gamma_X=\Gamma_Yf_*^G.
\tag{5.2.2}
\]

\SourceFootnote{10}{\LayerTarget{tohoku:footnote:10}
More generally, if $f$ is a continuous map from a space $X$ to a space
$Y$, then for sheaves $A$ on $X$ there are functorial homomorphisms
$f^{-1}(f_*(A))\to A$, from which the monomorphism (5.1.2) follows
immediately by virtue of the injection $A^G\to f_*(A)$.}

\TohokuContinuation{200}
